\section{Appendix B Plan: Scaled Hydrodynamic Equations}
\label{app:scaled-equations-plan}

\sourceeqrange{Appendix B Eqs. (B1)--(B6)}

\subsection*{Purpose and Scope}

This appendix gives an original guided derivation of the scaled
hydrodynamic equations.  The first pass derived the scale definitions and
dimensionless variables needed for Appendix B.  The present pass adds the
scaled continuity and momentum bookkeeping and treats the pressure tensor in
two branches: the expression printed in Eq.~(B3), and the
dimensional-source-consistent branch obtained by rescaling the pressure tensor
from Eqs.~(2.47)--(2.48).  The companion keeps these branches separate and
does not infer which branch, if either, was used in unpublished simulation
code.

\subsection*{Dimensional Objects}

The dimensional fields entering the scaling are:
\[
  n,\quad T,\quad v,\quad x,\quad t,\quad e_T,\quad \Pi,\quad \sigma .
\]
Here \(n\) is number density, \(T\) is temperature, \(v\) is velocity, and
\(e_T\), \(\Pi\), and \(\sigma\) have energy-density or stress dimensions.
The molecular volume parameter is \(v_0\), the particle mass is \(m\), the
van der Waals energy scale is \(\epsilon\), and \(k_B\) is Boltzmann's
constant.

\subsection*{Source Scale Choices}

The Section III.A simulation specialization used before Appendix B sets
\[
  K=0,\qquad C=\hbox{constant},
\]
and introduces the interface-length scale
\begin{equation}
  \ell=\left(\frac{C}{2k_Bv_0}\right)^{1/2}.\sourceeqbadge{B1}
  \label{eq:appb-ell-scale}
\end{equation}
The time scale is
\begin{equation}
  t_0=\frac{\ell^2}{\nu_0},
  \label{eq:appb-time-scale}
\end{equation}
where \(\nu_0\) is the source kinematic-viscosity scale.  The resulting
velocity scale is \(\ell/t_0\).  These are source scale choices, not derived
simulation results.

\subsection*{Scaled Variables in Eq. (B1)}

Appendix B then defines the dimensionless density, temperature, and velocity
variables
\begin{equation}
  \phi=v_0 n,\qquad
  \Theta=\frac{k_B T}{\epsilon},\qquad
  V=\frac{t_0}{\ell}v .\sourceeqbadge{B1}
  \label{eq:appb-scaled-variables}
\end{equation}
The spatial and time coordinates are also read in scaled units:
\[
  x_{\mathrm{sc}}=x/\ell,\qquad t_{\mathrm{sc}}=t/t_0 .
\]
The source suppresses separate symbols for the scaled coordinates after this
change.  The companion keeps the subscript here only to make the type change
visible.  The differential operators follow from the chain rule:
\begin{equation}
  \partial_t=t_0^{-1}\partial_{t_{\rm sc}},\qquad
  \partial_{x_i}=\ell^{-1}\partial_{i,\rm sc},\qquad
  \nabla_x=\ell^{-1}\nabla_{\rm sc},\qquad
  \Delta_x=\ell^{-2}\Delta_{\rm sc}.
  \label{eq:appb-operator-scales}
\end{equation}
For a vector flux, the same coordinate scaling gives
\(\nabla_x\cdot F=\ell^{-1}\nabla_{\rm sc}\cdot F\) before any field-specific
amplitude scale is substituted.  These operator identities are the only
source of powers of \(\ell\); they cannot create or remove a factor of
\(\phi\).

The substitution ledger used throughout Appendix B is:
\begin{center}
\begin{tabular}{@{}lll@{}}
\toprule
Dimensional object & Scaled substitution & Derivative consequence \\
\midrule
\(n\) & \(\phi/v_0\) & \(\nabla_x n=(v_0\ell)^{-1}\nabla\phi\) \\
\(T\) & \(\epsilon\Theta/k_B\) & \(\nabla_xT=\epsilon(k_B\ell)^{-1}\nabla\Theta\) \\
\(v\) & \((\ell/t_0)V\) & \(\partial_t v=(\ell/t_0^2)\partial_tV\) \\
\(x\) & \(\ell x_{\rm sc}\) & \(\nabla_x=\ell^{-1}\nabla_{\rm sc}\) \\
\(M\) & \(CT=2\epsilon v_0\ell^2\Theta\) & valid only for \(K=0\), constant \(C\) \\
\bottomrule
\end{tabular}
\end{center}

\subsection*{Stress, Energy, Gravity, and Transport Scales}

The natural pressure and stress scale is energy per molecular volume:
\[
  \Pi_{\mathrm{scale}}=\epsilon/v_0 .
\]
The same scale converts total energy density to the dimensionless quantity
used later in Appendix B:
\[
  E_T=\frac{v_0}{\epsilon}e_T .
\]
The dimensionless inertial parameter and gravity parameter are
\begin{equation}
  \sigma=\frac{\nu_0^2m}{\epsilon\ell^2}
  =
  \frac{m\ell^2}{\epsilon t_0^2},
  \qquad
  g^*=\frac{mg\ell}{\epsilon}.
  \label{eq:appb-sigma-gravity}
\end{equation}
The equality of the two forms of \(\sigma\) follows from
\eqref{eq:appb-time-scale}.  It is an algebraic scale identity.  The value of
\(\sigma\) used in a simulation is a modeling choice.

The source transport-coefficient branch used in Section III.A is
\[
  \eta=\zeta=\nu_0mn,\qquad
  \lambda=k_B\nu_0 n .
\]
These relations define the branch used for the later scaled equations.  The
momentum equation below uses this branch only for the viscous part of the
scaled stress.

\subsection*{Scaled Continuity and Momentum Bookkeeping}

From the number-density balance in Eq.~(2.35), the substitutions
\[
  n=\phi/v_0,\qquad v=(\ell/t_0)V,\qquad
  x=\ell x_{\mathrm{sc}},\qquad t=t_0t_{\mathrm{sc}}
\]
give
\[
  \partial_t n=\frac{1}{v_0t_0}\partial_{t_{\rm sc}}\phi,
  \qquad
  \nabla_x\cdot(nv)=\frac{1}{v_0t_0}\nabla_{\rm sc}\cdot(\phi V).
\]
Multiplying by \(v_0t_0\) gives the scaled local balance
\begin{equation}
  \frac{\partial\phi}{\partial t}
  =-\nabla\cdot(\phi V),\sourceeqbadge{B2}
  \label{eq:appb-scaled-continuity}
\end{equation}
where the derivatives in this appendix are with respect to the scaled
coordinates after the change of variables.

For the momentum balance, write \(\rho=mn\), scale the stress by
\(\epsilon/v_0\), and multiply the dimensional equation by
\(\ell v_0/\epsilon\).  A representative inertial term scales as
\[
  \frac{\ell v_0}{\epsilon}\partial_t(\rho v)
  =
  \frac{m\ell^2}{\epsilon t_0^2}\partial_{t_{\rm sc}}(\phi V),
\]
so the inertial coefficient becomes
\[
  \frac{m\ell^2}{\epsilon t_0^2}=\sigma.
\]
and the gravity coefficient becomes \(g^*=mg\ell/\epsilon\).  The same
substitutions give the term-by-term ledger
\begin{center}
\begin{tabular}{@{}p{0.35\linewidth}p{0.50\linewidth}@{}}
\toprule
Dimensional term & After multiplying by \(\ell v_0/\epsilon\) \\
\midrule
\(\partial_t(\rho v)\) &
\(\sigma\,\partial_{t_{\rm sc}}(\phi V)\) \\
\(\nabla_x\cdot(\rho vv)\) &
\(\sigma\,\nabla_{\rm sc}\cdot(\phi VV)\) \\
\(\nabla_x\cdot\Pi\), with \(\Pi=(\epsilon/v_0)\mathsf{M}\) &
\(\nabla_{\rm sc}\cdot\mathsf{M}\) \\
\(\rho g e_z\) &
\(g^*\phi e_z\) \\
\bottomrule
\end{tabular}
\end{center}
Using \eqref{eq:appb-scaled-continuity} to write the left side in
conservative form gives the scaled momentum pattern
\begin{equation}
  \sigma\left[
  \frac{\partial}{\partial t}(\phi V)
  +\nabla\cdot(\phi VV)\right]
  =
  -\nabla\cdot\mathsf{M}-g^*\phi e_z.
  \label{eq:appb-scaled-momentum-pattern}
\end{equation}
This is a scaling identity for the finite objects listed above.  It does not
by itself decide the pressure-tensor branch.

\subsection*{Branchwise Scaled Pressure Tensor}

Let
\[
  P_0(\phi,\Theta)=\frac{\Theta\phi}{1-\phi}-\phi^2 .
\]
Under the Appendix-B specialization \(K=0\) and constant \(C\), the coefficient
relation \(M=CT+K\) becomes \(M=CT\).  The source length scale
\(\ell^2=C/(2k_Bv_0)\) then converts the dimensional pressure tensor in
Eqs.~(2.47)--(2.48) into the dimensional-source-consistent scaled branch
\begin{equation}
  \mathsf{M}^{\mathrm{src}}_{ij}
  =
  \left[
    P_0-\Theta|\nabla\phi|^2
    -2\Theta\phi\,\Delta\phi
  \right]\delta_{ij}
  +2\Theta(\partial_i\phi)(\partial_j\phi)
  -\sigma\phi(\partial_iV_j+\partial_jV_i).
  \companionbranchbadge{scaled from 2.47--2.48}
  \label{eq:appb-source-consistent-b3}
\end{equation}
The factor \(\phi\) in the diagonal Laplacian term comes from the dimensional
term \(M n\Delta n\) in Eq.~(2.48), after \(n=\phi/v_0\) is substituted.
The common coefficient checks behind \eqref{eq:appb-source-consistent-b3} are
\begin{align}
  \frac{v_0}{\epsilon}p &= P_0,\label{eq:appb-b3-bulk-pressure-scale}\\
  \frac{v_0}{\epsilon}
    \left(-\frac{M}{2}|\nabla_x n|^2\right)
    &=-\Theta|\nabla\phi|^2,\label{eq:appb-b3-grad-square-scale}\\
  \frac{v_0}{\epsilon}(-Mn\Delta_x n)
    &=-2\Theta\phi\Delta\phi,\label{eq:appb-b3-laplacian-scale}\\
  \frac{v_0}{\epsilon}M(\partial_{x_i}n)(\partial_{x_j}n)
    &=2\Theta(\partial_i\phi)(\partial_j\phi).
    \label{eq:appb-b3-dyad-scale}
\end{align}
Here derivatives on \(\phi\) are scaled-coordinate derivatives.  The pressure
scale, gradient-square scale, and dyadic scale are common to both B3 branches;
the diagonal Laplacian row is the only row that separates the printed branch
from the dimensional-source branch.

The expression printed as Eq.~(B3) is kept as its own branch:
\begin{equation}
  \mathsf{M}^{\mathrm{print}}_{ij}
  =
  \left[
    P_0-\Theta|\nabla\phi|^2
    -2\Theta\,\Delta\phi
  \right]\delta_{ij}
  +2\Theta(\partial_i\phi)(\partial_j\phi)
  -\sigma\phi(\partial_iV_j+\partial_jV_i).\sourceeqbadge{B3 printed branch}
  \label{eq:appb-printed-b3-branch}
\end{equation}
The two branches differ only in the diagonal Laplacian factor.  The remaining
terms are common: the local bulk pressure \(P_0\), the gradient-square diagonal
term, the off-diagonal dyad, and the viscous scaled stress.  The branch
difference is therefore
\begin{equation}
  \mathsf{M}^{\mathrm{print}}_{ij}
  -\mathsf{M}^{\mathrm{src}}_{ij}
  =
  2\Theta(\phi-1)\Delta\phi\,\delta_{ij}.
  \label{eq:appb-b3-branch-residual}
\end{equation}
This expected nonzero residual is a source-branch marker.  Under the
dimensional-source derivation, the missing factor is a high-confidence
typographical-omission candidate in the printed formula.  This is not an
official erratum, not a silent correction of the source file, and not evidence
about the unpublished numerical implementation.

\begin{figure}[htbp]
\centering
\resizebox{0.88\linewidth}{!}{\begin{tikzpicture}[
  font=\small,
  node distance=8mm and 10mm,
  box/.style={draw=black, line width=0.45pt, rounded corners=2pt,
    align=center, text width=34mm, minimum height=11mm, fill=gray!8},
  branch/.style={draw=black, line width=0.45pt, rounded corners=2pt,
    align=center, text width=38mm, minimum height=13mm, fill=gray!14},
  open/.style={draw=black, dashed, line width=0.45pt, rounded corners=2pt,
    align=center, text width=40mm, minimum height=12mm, fill=white},
  arrow/.style={-{Latex[length=2.2mm]}, line width=0.5pt},
  darrow/.style={-{Latex[length=2.2mm]}, dashed, line width=0.5pt}
]
\node[box] (source) {Dimensional pressure tensor\\Eqs. (2.47)--(2.48)};
\node[box, right=of source] (scale) {Appendix-B scaling\\\(K=0,\ C\) constant};
\node[branch, below left=of scale] (dim) {Dimensional-source branch\\\(-2\Theta\phi\Delta\phi\)};
\node[branch, below right=of scale] (print) {Printed B3 branch\\\(-2\Theta\Delta\phi\)};
\node[box, below=of scale] (resid) {Branch residual\\\(2\Theta(\phi-1)\Delta\phi\,\delta_{ij}\)};
\node[open, below=of resid] (code) {Unpublished simulation-code branch\\not inferred};

\draw[arrow] (source) -- (scale);
\draw[arrow] (scale) -- (dim);
\draw[arrow] (scale) -- (print);
\draw[arrow] (dim) -- (resid);
\draw[arrow] (print) -- (resid);
\draw[darrow] (resid) -- (code);
\end{tikzpicture}
}
\caption{Appendix-B B3 branch map.  The dimensional-source scaling and the
printed Eq.~(B3) expression are kept as separate source branches.  The dashed
arrow records that the unpublished simulation-code branch is not inferred from
either printed branch.}
\label{fig:appendix-b-b3-branch-map}
\end{figure}

\subsection*{Direct Check of the B3 Diagonal Laplacian Factor}

The diagonal branch can be checked directly from the Appendix-B scale choices.
Use
\[
  \phi=v_0n,
  \qquad
  \Theta=\frac{k_BT}{\epsilon},
  \qquad
  x=\ell x_{\rm sc},
  \qquad
  \ell^2=\frac{C}{2k_Bv_0},
\]
together with the Section III.A specialization
\[
  K=0,
  \qquad C=\hbox{constant},
  \qquad M=CT .
\]
Then
\[
  M=CT=\frac{C\epsilon}{k_B}\Theta
    =2\epsilon v_0\ell^2\Theta,
  \qquad
  n=\frac{\phi}{v_0}.
\]
Because the derivatives in Eq.~\eqref{eq:appb-source-consistent-b3} are scaled
derivatives after the change of variables,
\[
  \nabla_x n=\frac{1}{v_0\ell}\nabla_{\rm sc}\phi,
  \qquad
  \Delta_x n=\frac{1}{v_0\ell^2}\Delta_{\rm sc}\phi .
\]
Scaling the dimensional pressure tensor by \(v_0/\epsilon\) gives
\begin{align}
  \frac{v_0}{\epsilon}\left(-Mn\Delta_x n\right)
  &=-\frac{M}{\epsilon v_0\ell^2}\phi\Delta_{\rm sc}\phi
    =-2\Theta\phi\Delta_{\rm sc}\phi,\label{eq:appb-direct-lap-factor}\\
  \frac{v_0}{\epsilon}\left(-\frac{M}{2}|\nabla_x n|^2\right)
  &=-\Theta|\nabla_{\rm sc}\phi|^2,\label{eq:appb-direct-grad-square}\\
  \frac{v_0}{\epsilon}M(\partial_{x_i}n)(\partial_{x_j}n)
  &=2\Theta(\partial_i\phi)(\partial_j\phi).
    \label{eq:appb-direct-gradient-stress}
\end{align}
Thus the dimensional-source branch formed from Eqs.~(2.47)--(2.48) and the
Appendix-B scalings produces the diagonal Laplacian term
\(-2\Theta\phi\Delta\phi\), not \(-2\Theta\Delta\phi\).  The factor
\(\phi\) enters only through the dimensional multiplier \(n\) in
\(n\Delta n\); coordinate scaling supplies powers of \(\ell\), and the
pressure scale supplies \(v_0/\epsilon\), but neither supplies nor cancels a
density factor.  The printed branch therefore differs by
\[
  \mathsf{M}^{\rm print}_{ij}-\mathsf{M}^{\rm src}_{ij}
  =2\Theta(\phi-1)\Delta\phi\,\delta_{ij}.
\]
The corresponding scalar diagonal contribution to the momentum equation is
\[
  -\partial_j\left(\mathsf{M}^{\rm print}_{ij}
       -\mathsf{M}^{\rm src}_{ij}\right)
  =-\partial_i\left[2\Theta(\phi-1)\Delta\phi\right],
\]
which is generically nonzero.  For example, with constant \(\Theta\) and the
one-dimensional test field \(\phi(x)=1+x^2\), this force residual is
\(-8\Theta x\).

This direct scaling check rules out the common algebraic escape routes: the
coordinate scaling, stress scale, definition \(\phi=v_0n\), and the
\(K=0\), constant-\(C\) specialization all retain the factor \(\phi\).
Removing the factor would require replacing the dimensional source term
\(n\Delta n\) by \(\Delta n\), which is not the dimensional pressure-tensor term
being scaled.  The result makes the missing factor a high-confidence
typographical-omission candidate in the printed Eq.~(B3) branch under the
dimensional-source derivation.  It is still not an official erratum, not a
statement of author intent, and not evidence about the unpublished
simulation-code branch.

\subsection*{Scaled Total-Energy Density}

Under the same \(K=0\), constant-\(C\) Appendix-B specialization, the source
scaled total-energy density in Eq.~(B4) is
\begin{equation}
  E_T=\phi\Theta-\phi^2+\frac{\sigma}{2}\phi |V|^2 .\sourceeqbadge{B4}
  \label{eq:appb-scaled-energy-density}
\end{equation}
The first two terms are the local van der Waals internal-energy branch after
multiplication by \(v_0/\epsilon\).  The kinetic term follows from
\[
  \frac{v_0}{\epsilon}\frac{mn}{2}|v|^2
  =
  \frac{1}{2}\frac{m\ell^2}{\epsilon t_0^2}\phi |V|^2
  =
  \frac{\sigma}{2}\phi |V|^2 .
\]
This is a local density scaling.  It is not a global energy theorem and does
not include a decision about wall heat or surface-energy transfer.

\subsection*{Scaled Total-Energy Balance}

Scaling the local total-energy balance in Eq.~(2.39) by
\(v_0t_0/\epsilon\), using \(t_0=\ell^2/\nu_0\), and substituting the
transport branch \(\lambda=k_B\nu_0n\), gives the Appendix-B energy-balance
pattern.  The heat term is the cleanest coefficient check.  First scale the heat flux
inside the divergence:
\[
  \lambda\nabla_xT
  =k_B\nu_0\frac{\phi}{v_0}
    \frac{\epsilon}{k_B\ell}\nabla_{\rm sc}\Theta
  =\frac{\epsilon\nu_0}{v_0\ell}\phi\nabla_{\rm sc}\Theta .
\]
Then include the divergence scale and the balance prefactor:
\begin{align}
  \frac{v_0t_0}{\epsilon}\nabla_x\cdot(\lambda\nabla_xT)
  &=\frac{v_0t_0}{\epsilon\ell^2}
    \nabla_{\rm sc}\cdot\left(k_B\nu_0\frac{\phi}{v_0}
    \frac{\epsilon}{k_B}\nabla_{\rm sc}\Theta\right) \\
  &=\frac{t_0\nu_0}{\ell^2}\nabla_{\rm sc}\cdot(\phi\nabla_{\rm sc}\Theta)
   =\nabla_{\rm sc}\cdot(\phi\nabla_{\rm sc}\Theta).
  \label{eq:appb-heat-term-scaling-step}
\end{align}
The gravity-power term similarly becomes
\[
  \frac{v_0t_0}{\epsilon}\rho g v_z
  =\frac{mg\ell}{\epsilon}\phi V_z
  =g^*\phi V_z .
\]
Thus the scaled balance is
\begin{equation}
  \frac{\partial E_T}{\partial t}
  +\nabla\cdot\left(E_T V+\mathsf{M}\cdot V\right)
  =
  \nabla\cdot(\phi\nabla\Theta)-g^*\phi V_z .\sourceeqbadge{B5}
  \label{eq:appb-scaled-energy-balance}
\end{equation}
Here \(\mathsf{M}\) is the scaled stress tensor from the preceding branchwise
discussion.  The scaling of \(\nabla\cdot(\lambda\nabla T)\) gives
\(\nabla\cdot(\phi\nabla\Theta)\); the scaling of gravitational power gives
\(-g^*\phi V_z\).  Choosing between the printed and dimensional-source B3
branches is not needed for these scale factors, but any evaluation of
\(\mathsf{M}\cdot V\) must state which branch is being used.

\subsection*{Scaled Density Boundary Condition}

The density boundary condition used for the simulations is stated before
Appendix B in Eq.~(3.2).  Its dimensional form can be read as
\[
  2\sqrt{2}\,\ell\,\nu\cdot\nabla_x n=n-\frac{5}{2}n_c.
\]
With \(n=\phi/v_0\), scaled normal derivative
\(\nu\cdot\nabla_{\mathrm{sc}}\phi=v_0\ell\,\nu\cdot\nabla n\), and the
van der Waals critical density \(v_0n_c=1/3\), the two sides scale as
\begin{align}
  2\sqrt{2}\,\ell\,\nu\cdot\nabla_x n
    &=\frac{2\sqrt{2}}{v_0}\,\nu\cdot\nabla_{\rm sc}\phi,\\
  n-\frac{5}{2}n_c
    &=\frac{1}{v_0}\left(\phi-\frac{5}{6}\right).
\end{align}
Multiplying by \(v_0\) gives
\[
  2\sqrt{2}\,\nu\cdot\nabla_{\rm sc}\phi
  =\phi-\frac{5}{6}.
\]
Equivalently,
\begin{equation}
  \nu\cdot\nabla\phi=\frac{\phi-5/6}{2\sqrt{2}} .\sourceeqbadge{B6}
  \label{eq:appb-scaled-density-boundary}
\end{equation}
This is a density wall branch.  It does not by itself specify the heat
boundary data or the surface-energy transfer appearing in the global entropy
statement of Eq.~(2.53).

\subsection*{What Is Deferred}

The following are deliberately not derived in this pass:

\begin{itemize}
\item any numerical reproduction or sensitivity study for the Section III simulations;
\item global wall/surface-transfer closure beyond the local density boundary
  condition;
\item the unpublished simulation-code branch.
\end{itemize}

\subsection*{Verification Checks}

The companion checks the scale definitions with finite algebra and dimensions:
\[
  \phi,\quad \Theta,\quad V,\quad \sigma,\quad g^*,\quad
  \frac{v_0e_T}{\epsilon}
\]
are dimensionless, and
\[
  \frac{\nu_0^2m}{\epsilon\ell^2}
  =
  \frac{m\ell^2}{\epsilon t_0^2}
\]
after substituting \(t_0=\ell^2/\nu_0\).  These checks do not derive the
scaled energy equation or validate any numerical run.

The momentum/B3 pass additionally checks:
\[
  \mathsf{M}^{\mathrm{print}}_{ij}
  -\mathsf{M}^{\mathrm{src}}_{ij}
  =
  2\Theta(\phi-1)\Delta\phi\,\delta_{ij},
\]
the inertial coefficient \(\sigma\), the viscous coefficient
\(\sigma\phi\), direct scaling of the diagonal Laplacian factor from
\(-Mn\Delta n\), the one-dimensional force residual, and the negative control
that replacing \(n\Delta n\) by \(\Delta n\) is not source-grounded.  These are
finite algebraic checks only.

The energy/boundary pass checks:
\[
  E_T,\qquad \nabla\cdot(\phi\nabla\Theta),\qquad
  g^*\phi V_z,\qquad
  \nu\cdot\nabla\phi=\frac{\phi-5/6}{2\sqrt{2}} .
\]
It also keeps an expected nonzero marker for omitted global wall/surface
transfer.

The executable mirrors for the Appendix B derivations are:
\begin{center}
\small
\begin{tabular}{@{}l@{}}
\texttt{scripts/python/appendix\_b\_scaling\_definitions\_sympy.py}\\
\texttt{scripts/python/appendix\_b\_momentum\_b3\_sympy.py}\\
\texttt{scripts/python/appendix\_b\_b3\_typo\_candidate\_sympy.py}\\
\texttt{scripts/python/appendix\_b\_energy\_boundary\_sympy.py}\\
\texttt{scripts/wolfram/appendix\_b\_scaling\_definitions\_checks.wls}\\
\texttt{scripts/wolfram/appendix\_b\_momentum\_b3\_checks.wls}\\
\texttt{scripts/wolfram/appendix\_b\_b3\_typo\_candidate\_checks.wls}\\
\texttt{scripts/wolfram/appendix\_b\_energy\_boundary\_checks.wls}\\
\texttt{scripts/wolfram/appendix\_b\_transition\_trace.wls}\\
\texttt{tests/test\_appendix\_b\_scaling\_definitions.py}\\
\texttt{tests/test\_appendix\_b\_momentum\_b3.py}\\
\texttt{tests/test\_appendix\_b\_b3\_typo\_candidate.py}\\
\texttt{tests/test\_appendix\_b\_energy\_boundary.py}
\end{tabular}
\end{center}
These scripts mirror the displayed nondimensional substitutions, branch
residuals, dimensional checks, and boundary-condition scaling.  They do not
infer unpublished implementation behavior or reproduce the Section III
simulations.
The transition trace records the source object, substitution, intermediate
expression, expected scaled expression, and residual for Eqs.~(B1)--(B6).  Its
negative controls deliberately omit the coordinate factor, density factor, or
source gravity sign, or mix dimensional and scaled velocity symbols.

\subsection*{Source-Procedure Completion Status}

The Appendix-B rows in the exhaustive inventory are now covered branchwise:
\begin{itemize}
\item Row (B1) and the surrounding scaled-coordinate convention are
  handled by the substitution ledger and operator-scale identities.
\item The unnumbered scaled continuity equation and row (B2) follow from
  source rows (2.35)--(2.36) by the displayed
  nondimensionalization.
\item Row (B3) is closed only as a branchwise source-printing analysis:
  the dimensional-source scaling gives the \(-2\Theta\phi\Delta\phi\) term,
  while the printed expression lacks \(\phi\).  The companion classifies this
  as a high-confidence typographical-omission candidate under the
  dimensional-source derivation, not as an official erratum, author-intent
  claim, or simulation-code inference.
\item Rows (B4)--(B6) are covered by the total-energy density,
  total-energy balance, heat-flux scaling, gravity-power scaling, and density
  boundary-condition substitutions above.
\end{itemize}
No unpublished numerical implementation branch is inferred from these
calculations.

\subsection*{Tagged verification complement}

\OnukiLinkAppendixB{The tagged Appendix B complement} presents the
dimensional-to-scaled hand calculation, traces, and printed-branch negative
control. It does not state an official erratum or identify unpublished code.

\subsection*{Open issues}

The companion must not silently correct the printed source and must not infer
which branch was used by unpublished simulation code.  The Section III guide
will discuss the simulation scenarios as source reading, not as numerical
reproduction.
